% !TITLE 先秦
% !SUMMARY 一切的开始
% !ORDER 0

\begin{intro}
我们的旅程，从先秦开始，这是一切的开始、一切的起点。

三千年前的黄河之畔，一位君子爱上了一位淑女，却不知她的心意。他朝思暮想辗转反侧，终于在一个夜晚写下了一首诗。这就是后来的《关雎》。那个夜晚，文学诞生了。当然，中华文学的起源在此之前就已出现，在甲骨文的卜辞里，在田亩间的歌唱里，在宗庙祭祀的乐章里，在大胜而归的篝火晚会上，但当第一篇诗出现在大地上时，文学终于以一种独立的姿态展现在我们面前。此后，风骚雅颂，秦文汉赋，唐诗宋词，元曲明清小说，直到今天我们所看到的一切汉字、一切中文文学，就这样流淌了出来。

而回到最初，这条大河的源头，叫做《诗经》和《楚辞》。

《诗经》三百零五篇，收的是西周初年到春秋中叶五百年间的歌谣与乐章。它是北方的声音——是在黄河流域，史官文化的土壤里长出来的。风、雅、颂三种体例中，最动人的是十五国风，那些“饥者歌其食，劳者歌其事”的句子，没有作者署名，却比任何署名之作都更直接地触碰了人类共通的情感。文学的技巧在《诗经》里已经得到初步的使用——铺陈直叙的赋，以此喻彼的比，托物起兴的兴。赋比兴的写作技巧也成为了后世文学的重要技法，为什么歌唱祖国里，为什么要先唱“五星红旗迎风飘扬”，这就是起兴，这就是《诗经》的痕迹。

如果说《诗经》是北方的，是世俗的，那么《楚辞》就是南方的，是浪漫的。它是从长江流域，巫官文化的土壤里长出来的。它比《诗经》晚了几百年，气象却完全不同——如果说《诗经》是素白的棉麻，《楚辞》就是织金的锦缎。屈原的出现是一个标志：在此之前，诗歌是集体的、匿名的、口耳相传的；在此之后，有了一个人的名字，有了一个个体的全部情感与命运。于是文学不再只是集体的声音，也能承载个人的表达。那些在历史里姓名模糊的人物，借助文学留下了清晰的身影。今天的我们已经很难说出唐朝都有哪些皇帝，都做了什么说了什么，可我们都记得李白曾经在一个夜晚抬起头对着天上的明月思念故乡。这就是文学的价值。

《诗经》的现实主义与《楚辞》的浪漫主义，北方的质朴与南方的瑰丽，风与骚——这两条河流，从先秦大地的东西两端发源，汇成了此后两千年中国文学的干流。让我们踏入这条河流吧！
\end{intro}
